\subsection{LENGTH AND REDUCED DECOMPOSITIONS}

DEFINITION 1. Let \(w \in W\) . The length of \(w\) (with respect to \(S\) ), denoted by \(l_{S}(w)\) or simply by \(l(w)\) , is the smallest integer \(q \geqslant 0\) such that \(w\) is the product of a sequence of \(q\) elements of \(S\) . A reduced decomposition of \(w\) (with respect to \(S\) ) is any sequence \(s = (s_1, \ldots, s_q)\) of elements of \(S\) such that \(w = s_1 \ldots s_q\) and \(q = l(w)\) .

Thus 1 is the unique element of length 0 and \(S\) consists of the elements of length 1.

PROPOSITION 1. Let w and w' be in W. We have the formulas:

\[
l (w w ^ {\prime}) \leqslant l (w) + l (w ^ {\prime}),\tag{1}
\]

\[
l (w ^ {- 1}) = l (w),\tag{2}
\]

\[
\left| l (w) - l \left(w ^ {\prime}\right) \right| \leqslant l \left(w w ^ {\prime - 1}\right).\tag{3}
\]

Let \((s_{1},\ldots,s_{p})\) and \((s_{1}^{\prime},\ldots,s_{q}^{\prime})\) be reduced decompositions of w and \(w^{\prime}\) respectively. Thus \(l(w)=p\) and \(l(w^{\prime})=q\) . Since \(ww^{\prime}=s_{1}\ldots s_{p}s_{1}^{\prime}\ldots s_{q}^{\prime}\) , we have \(l(ww^{\prime})\leqslant p+q\) , proving (1). Since \(S=S^{-1}\) and \(w^{-1}=s_{p}^{-1}\ldots s_{1}^{-1}\) , we have \(l(w^{-1})\leqslant p=l(w)\) . Replacing w by \(w^{-1}\) gives the opposite inequality \(l(w)\leqslant l(w^{-1})\) , proving (2). Replacing w by \(ww^{\prime-1}\) in (1) and (2) gives the relations

\[
l (w) - l \left(w ^ {\prime}\right) \leqslant l \left(w w ^ {\prime - 1}\right),\tag{4}
\]

\[
l \left(w w ^ {\prime - 1}\right) = l \left(w ^ {\prime} w ^ {- 1}\right).\tag{5}
\]

Exchanging \(w\) and \(w'\) in (4) gives \(l(w') - l(w) \leqslant l(ww'^{-1})\) by (5), proving (3).

COROLLARY. Let \(\mathbf{s} = (s_{1}, \ldots, s_{p})\) and \(\mathbf{s}' = (s_{1}', \ldots, s_{q}')\) be two sequences of elements of S such that \(w = s_{1} \ldots s_{p}\) and \(w' = s_{1}' \ldots s_{q}'\) . If the sequence \((s_{1}, \ldots, s_{p}, s_{1}', \ldots, s_{q}')\) is a reduced decomposition of \(ww'\) , then s is a reduced decomposition of w and \(s'\) is one of \(w'\) .

By hypothesis, \(l(w) \leqslant p\) , \(l(w') \leqslant q\) and \(l(ww') = p + q\) . By (1), we must have \(l(w) = p\) and \(l(w') = q\) , hence the corollary.

Remark In view of formulas (1) and (2), the formula \(d(w, w') = l(ww'^{-1})\) defines a distance on W, invariant under right translations.

